\section{Method}\label{sec:method}

The method works in two stages. The first stage scores every equation of the database on its own
and keeps the best as candidates. The second stage reranks the candidates with a learned model that
compares each candidate with other equations.

\subsection{Stage 1: candidate retrieval}\label{sec:stage1}

The first stage keeps the $N_1$ equations with the highest score as the candidate set $C_q$. The
score mixes one text term and one variable term,
\begin{equation}
  s_1(q, e_j) = w_{\mathrm{t}} \cos\!\big(\mathrm{tfidf}(c), \mathrm{tfidf}(t_j)\big)
              + w_{\mathrm{v}} \, J\!\big(\mathcal{Q}, V(e_j)\big) ,
  \label{eq:stage1}
\end{equation}
where $t_j$ joins the context text, the expression, and the domain label of equation $e_j$. The
first term, the text similarity, compares the description with this text through
term-frequency--inverse-document-frequency (TF-IDF) vectors \citep{Salton} over word unigrams and
bigrams. The second term, the variable overlap, is the Jaccard similarity
$J(A, B) = |A \cap B| / |A \cup B|$ between the query variables and the variables of the equation.
Ranking the whole database by Eq.~\eqref{eq:stage1} is a classical information-retrieval method
\citep{Manning}, and Section~\ref{sec:setup} uses it, without a learned model or a candidate set,
as the baseline.

The second stage reranks only the candidates, so the fraction of ground-truth equations among them,
the candidate coverage, bounds what the second stage can achieve. We set $N_1$ so large that every
ground-truth equation enters the candidates (Section~\ref{sec:compared}).

\subsection{Stage 2: set-aware reranking}\label{sec:stage2}

The second stage judges each candidate against the equations already at hand.
Table~\ref{tab:features} lists its ten features. Seven compare the query with a single equation:
two measure the topical closeness of the description ($f_1$, $f_3$), one matches the domain label
against the description ($f_7$), and four measure how the variables of the equation relate to the
query variables ($f_2$, $f_4$, $f_5$, $f_6$). The other three are the set-aware features that this
study introduces. They compare the candidate with a reference set $R$: the top $m$ equations of
Stage~1 in static reranking, and the equations already selected in sequential selection
(Section~\ref{sec:seq}). A candidate that belongs to $R$ is compared with $R$ without itself.

Complementarity $g_{\mathrm{Comp}}$ asks what the candidate adds: the fraction of query variables
that appear in the candidate but in no equation of the reference set. Coherence $g_{\mathrm{Coh}}$
asks whether the candidate connects to the reference set: the fraction of its variables shared
with the set. Domain agreement $g_{\mathrm{Dom}}$ asks whether the candidate comes from the same
topic: the fraction of reference equations whose domain label equals the candidate's. The first
two encode how equations determine one another's variables, and the third encodes topical
consistency.

\begin{table}[t]
\centering
\caption{Features of the reranker. $\mathcal{Q} = \mathcal{I} \cup \mathcal{O}$ is the set of query
variables, $t_j$ the text of equation $e_j$, $R$ the reference set, and
$V(R) = \bigcup_{e \in R} V(e)$. The group marks whether a feature reads variable overlap or
topical closeness.}
\label{tab:features}
\small
\begin{tabular}{llll}
\toprule
 & Feature & Definition & Group \\
\midrule
$f_1$ & Text similarity     & $\cos(\mathrm{tfidf}(c), \mathrm{tfidf}(t_j))$ & topic \\
$f_2$ & I/O Jaccard         & $J(\mathcal{Q}, V(e_j))$ & variable \\
$f_3$ & Latent similarity   & $\cos(\mathrm{svd}(c, \mathcal{Q}), \mathrm{svd}(t_j))$ & topic \\
$f_4$ & Input coverage      & $|\mathcal{I} \cap V(e_j)| / |\mathcal{I}|$ & variable \\
$f_5$ & Output coverage     & $|\mathcal{O} \cap V(e_j)| / |\mathcal{O}|$ & variable \\
$f_6$ & Specificity         & $|\mathcal{Q} \cap V(e_j)| / |V(e_j)|$ & variable \\
$f_7$ & Domain match        & $1$ if $\mathrm{dom}(e_j)$ occurs in $c$ or $c$ in $\mathrm{dom}(e_j)$ & topic \\
\midrule
$g_{\mathrm{Comp}}$ & Complementarity  & $|(V(e_j) \cap \mathcal{Q}) \setminus V(R)| / |\mathcal{Q}|$ & variable \\
$g_{\mathrm{Coh}}$  & Coherence        & $|V(e_j) \cap V(R)| / |V(e_j)|$ & variable \\
$g_{\mathrm{Dom}}$  & Domain agreement & $|\{e \in R : \mathrm{dom}(e) = \mathrm{dom}(e_j)\}| / |R|$ & topic \\
\bottomrule
\end{tabular}
\end{table}

The latent similarity $f_3$ applies latent semantic analysis \citep{LSA}: a truncated
singular-value decomposition factorizes the TF-IDF matrix of all equation texts into 256
dimensions, the query is projected into the same space, and $f_3$ is the cosine similarity there.
The similarity grows when two texts contain words that appear in similar texts, even without shared
words. Unlike $f_1$, the query text of $f_3$ also lists the input and output variable symbols.

A multilayer perceptron (MLP) with two hidden layers of \Hidden{} units maps the ten features
$x(q, e)$ of a candidate to a score $f(x(q, e))$. We train it with a pairwise ranking loss
\citep{LTR} that asks every correct candidate to outscore the incorrect candidates of the same case
by at least a margin $\gamma$. For a training case $q$, a ground-truth equation
$e^{+} \in \mathcal{R}_q \cap C_q$, and an incorrect candidate
$e^{-} \in C_q \setminus \mathcal{R}_q$,
\begin{equation}
  \ell(q, e^{+}, e^{-}) = \max\!\big(0,\; \gamma - f(x(q, e^{+})) + f(x(q, e^{-}))\big) .
  \label{eq:loss}
\end{equation}
For each training case we choose $n_{\mathrm{neg}}$ incorrect candidates and pair every correct
candidate with each of them. We compare two ways of choosing them and let the validation cases
decide (Section~\ref{sec:compared}): at random, or the incorrect candidates that the MLP scores
highest at that point of training (hard negatives). With a large $N_1$, most randomly chosen
incorrect candidates already trail the correct ones by more than the margin and contribute no loss.
The hyperparameters of the method, the values fixed before training, are $N_1$, $m$, the
first-stage weights, $\gamma$, $n_{\mathrm{neg}}$, the number of epochs, and the learning rate
(\ref{app:hyper}).

\subsection{Sequential selection}\label{sec:seq}

Static reranking fixes the reference set to the top $m$ equations of Stage~1 and scores all
candidates in one pass. When a case has many ground-truth equations, the equations that rise to
the top drift away from this fixed set. Sequential selection therefore replaces the reference set
with the equations selected so far and selects one equation at a time. It first takes the candidate
with the highest static score. It then recomputes the set-aware features of the remaining
candidates against the equations selected so far, takes the candidate with the highest score, and
repeats until it has selected 20 equations; the unselected candidates follow in the order of their
static scores.

Sequential selection comes in two variants that differ in how the MLP is trained. Greedy selection
at inference uses the MLP trained for static reranking. Learned greedy selection also trains on the
sequential contexts, with teacher forcing: it orders the ground-truth equations among the
candidates of a training case by their Stage-1 rank, $e^{+}_1, e^{+}_2, \dots$, treats
$e^{+}_1, \dots, e^{+}_{t-1}$ as the equations already selected at step $t$, and adds the loss of
Eq.~\eqref{eq:loss} for the next ground-truth equation $e^{+}_t$ with the features computed against
that reference set. A case contributes at most eight steps. Training thus asks the same question as
inference: which equation to select next, given the ones already selected.
